\chapter{Introduction}\label{ch:introduction}
This study asks how far repeated corrosion photographs can support an inference about ferrocement condition. Its central concern is the boundary between estimating an observable surface quantity and inferring a structural endpoint that has much less supervision.

\section{Motivation and problem}
Ferrocement combines a cementitious matrix with distributed steel-mesh reinforcement. For the laboratory specimens considered here, external rust evidence can be recorded repeatedly while structural testing is destructive. This asymmetry makes visual monitoring attractive but complicates the interpretation of an image-based predictor. An accurate estimate of visible corrosion is not automatically an accurate estimate of lost wire area, residual capacity, or future service life.

The physical specimens, exposure campaigns, and terminal experiments originate in predecessor work \citep{hossain_thesis,artiste2025}. The present activity concerns subsequent data preparation, image representations, supervised models, evaluation, and interpretation of the saved computational results. It does not claim authorship of the original experimental campaign or a fresh reproduction of its image-processing protocol.

A recurring scientific question is whether a model learns a transferable relationship or a feature of the experimental mixture. If each campaign uses its own mesh configuration and exposure conditions, a strong pooled capacity score can partly reflect the separation between campaigns. Testing on new specimens from that mixture and testing on an unseen campaign therefore ask materially different questions.

\section{Research questions}
The report organizes the evidence around five questions:
\begin{enumerate}[label=\textbf{RQ\arabic*.}]
\item How consistently can visible surface corrosion be estimated across distinct image representations, and how does label adjacency limit that interpretation?
\item How well do the available inputs predict the two terminal structural endpoints, and how stable are the results under stricter evaluation?
\item Do image-derived corrosion descriptors add predictive value beyond specimen and exposure metadata for terminal ultimate load?
\item How do grouped specimen evaluation, treatment exclusion, campaign exclusion, and mesh stratification change the defensible claim?
\item What does a model-derived degradation and threshold-crossing pipeline establish when measured structural trajectories and lifetime outcomes are absent?
\end{enumerate}
The prepared four-class classification branch is reported as a development status, not as a sixth answered performance question.

\section{Contributions and scope}
The first contribution is a synthesis that separates target observability from model sophistication. The same image collection supports a well-supervised surface task and a sparsely supervised structural task; these should not be assessed with the same evidential expectations.

The second contribution is a focused interpretation of the terminal-load ablation. Image-only, metadata-only, and combined models are compared under the same specimen partitions, with the scope of the prediction target made explicit. The experiment assesses terminal load from terminal held-out observations, even though training uses the broader image history with repeated terminal labels.

The third contribution is a verification and reporting framework for the saved evidence. It checks specimen disjointness, reproduces key metrics from saved predictions, resolves conflicting model-selection summaries, and documents figure provenance. This supports reproduction of the reported aggregation; it does not establish that the historically trained models can currently be regenerated without source repairs.

The conclusions are restricted to the available specimens, representations, and evaluation designs. They do not identify the independent causal effects of mesh, chloride, or duration. They do not validate a structural safety decision, a remaining-life estimator, or a four-class classifier. They also do not establish that surface images contain no structural information: the distinction is between a physical relationship that may exist and a stable predictive relationship demonstrated by these experiments.

\section{Organization}
Chapter~\ref{ch:background} positions the study in imaging and evaluation research. Chapters~\ref{ch:data} and \ref{ch:methods} define observations and methods. Chapters~\ref{ch:surface}--\ref{ch:proxy} answer the research questions using saved evidence. The discussion and limitations then connect those observations to concrete research decisions. Reproduction details are confined to Appendix~\ref{app:reproduction} so the scientific narrative remains readable.
